\scanpage{115}
% ms. line 01
Soit $f(z)=\prod_1^\infty(1+\lambda_iz)$ une fonction entière
% ms. line 02
de \uline{genre} $1$, posons $\alpha_i=|\lambda_i|$, $F(z)=\prod_{\unsure{1|i}}^\infty(1+\alpha_i z)$.
% ms. line 03
Alors \uline{$f$ et $F$ ont $=$ ordre} (qui s'exprime donc
% ms. line 04
aussi en fonction de la suite des $|\alpha_i|$). En effet,
% ms. line 05
l'ordre de $f\leq$ ordre de $F$ (évident). D'autre
% ms. line 06
part, ordre de $F=$ ordre de la suite des $|\alpha_i|$
% ms. line 07
(voir ci-dessus) $\leq$ ordre de $f$ (th. de Hadamard).
\begin{center}\rule{1.5cm}{0.4pt}\end{center}

% ms. line 08
On peut parler de l'ordre d'un ensemble \inserted{dénombrable} de
% ms. line 09
nombres complexes, \deleted{indépendamment de} \inserted{$=$ \unsure{non} \unsure{rangés|compris} en suite :}
% ms. line 10
on prend la borne inf. des $s>0$ tels que
% ms. line 11
$\sum|z_i|^s<+\infty$. Suite à décroissance rapide $=$
% ms. line 12
d'ordre zéro.

% ms. line 13
Ordre d'une suite dans un espace de Banach,
% ms. line 14
ou de Fréchet. (deux définitions, suivant qu'on
% ms. line 15
prend sommabilité forte, ou sommabilité absolue)

% ms. line 16
Ordre d'une partie bornée (\deleted{plus petit} \inserted{borne inférieure} des
% ms. line 17
$s>0$ tels que la partie soit contenue dans l'enveloppe
% ms. line 18
convexe cerclée fermée d'une suite d'ordre $s$)
% archive sheet number: 116
